\section{Graph-uniform and modified Tian identification: complete proofs}
\label{sec:annex-p16-graph-uniform-identification}

\begin{annexitem}{annex:p16.graph-uniform-equivalences}{thm:p16-graph-uniform-equivalences}{The quantified equivalence and the four-valued negative witness}
\begin{proof}
For any fixed finite nonempty alphabet family, use independent private
variables with positive masses. Declare all required directed arguments
as ignored, and add independent ignored pair-specific latent inputs
for the required bidirected edges. This gives the prescribed graph
and a compatible positive product law. Thus the positive model classes in the outer quantifier are
nonempty. Every branch condition in ID uses only vertex sets,
directed ancestry or district membership. Its positive support invariant
ensures that no numerical denominator produces an additional branch.
The success/failure decision consequently depends only on \(G,X,Y\).
One can determine it on any positive compatible law; the same decision
is reached on every other finite alphabet family and positive law.

If ID succeeds, \cref{thm:p16-id-success-soundness} identifies its
expression with the intervention kernel for every compatible model.
That theorem permits every finite nonempty observed domain, including
singletons. Its sums and conditional products form one finite formula
whose index sets are interpreted on the supplied alphabets. Equal
observed laws in one family give equal values of this expression.
This proves graph-uniform identification, with the outer alphabet
quantifier retained. The same successful trace gives a uniform finite
do-calculus derivation by \cref{thm:p16-id-do-utility}.

If ID fails, run the unchanged graph recursion on four values at
every vertex. There are no singleton coordinates to remove in this
family. The previously proved failure-extraction theorem constructs
a hedge for the original \(X,Y\), and
\cref{thm:p16-hedge-four-state-routing} gives two models on this
four-valued family with equal strictly positive full observed laws
but different original intervention kernels. These models violate
one instance of the universal alphabet quantifier. Therefore the
query is not graph-uniformly identified. This is sufficient to prove
that graph-uniform identification implies ID success; a countermodel
on every other alphabet family is not required by that quantifier.

The extraction theorem also proves that absence of an original-query
hedge implies success, since every failed run would supply one.
Conversely, existence of any hedge for the complete query supplies
the same kind of positive four-valued countermodels. Success soundness
would make their effects equal if ID succeeded, a contradiction.
Thus success is equivalent to absence of such a hedge.

A uniform finite do-calculus derivation with observed-law output is
sound in every positive finite model by the three-rule soundness
theorem and ordinary probability identities. Applying it to two
models in any one family with equal observed laws yields equal
effects. It therefore implies graph-uniform identification. Combined
with the derivation supplied by every successful ID trace, this
establishes all four equivalences.

Finally the earlier singleton theorem concerns one specified family
and proves the exact effect-preserving reduction for it. A variable
constant in that family need not be constant in the four-valued
family used above. Removing it from the universally quantified
graph would change the outer question. Keeping the two quantifiers
separate proves the final assertion of the statement.
\end{proof}
\end{annexitem}

\begin{annexitem}{annex:p16.graph-uniform-all-effects}{thm:p16-graph-uniform-all-effects}{The graph-level criterion across all finite alphabets}
\begin{proof}
Suppose no directed edge has its endpoints in one maximal district.
If some ID query failed, its failed random graph would be one
district and would contain a directed path from a local action
vertex to its response. This path contains an edge whose endpoints
lie in that failed district, and hence in one original maximal
district, a contradiction. This is the graph argument already
proved in \cref{thm:p16-hedge-all-effects}; it is independent
of alphabet sizes. Every query therefore succeeds, and the preceding
uniform equivalence supplies its formula and do-calculus derivation
on every finite nonempty alphabet family.

Conversely a directed edge \(a\to b\) within one district has
a simple bidirected path between its endpoints. Keeping that path
and the single directed edge gives a forest with roots all its
vertices except \(a\); deleting \(a\) gives a smaller forest
with the same roots. As proved in \cref{thm:p16-hedge-all-effects},
this is a hedge for action \(\{a\}\) and full nonaction
response \(N\setminus\{a\}\). It has positive binary
countermodels, or the available four-valued countermodels. Either
one alphabet family disproves uniform identification of that query,
so not every query is uniformly identified. Bidirected connectivity
is exactly district membership, giving the equivalent path wording.
The empty graph and empty queries are included by the unit kernel.
\end{proof}
\end{annexitem}

\begin{annexitem}{annex:p16.graph-uniform-action-scope}{thm:p16-graph-uniform-action-scope}{A positive counterexample to enlarging the action set of an obstruction}
\begin{proof}
In the stated two models, observational substitution gives \(A=U\)
and \(M=E\) in both. The equation for \(Y\) is the same.
For each \(a\), the probability of \((A=a,M=m,Y=y)\) is
\[
 \tfrac12\,q(m)\,q(y\oplus m),\qquad q(0)=3/4,\quad q(1)=1/4.
\]
At \((m,y)=(0,0),(0,1),(1,0),(1,1)\), these products are
\(9/32,3/32,1/32,3/32\), respectively. They are positive
and sum to \(1/2\) for each \(a\), proving equality and
normalization of the full observed laws.

After \(A=a\), model 1 gives
\(M=a\oplus U\oplus E\), a fair bit, and hence a fair
\(Y=M\oplus E_Y\). Model 2 gives \(Y=E\oplus E_Y\),
whose one-probability is
\((1/4)(3/4)+(3/4)(1/4)=3/8\). Thus the action \(\{A\}\)
does not uniformly identify the response \(Y\).
The bow on \(A,M\) and its singleton smaller forest \(M\)
also form a hedge for this query: the root \(M\) is an ancestor
of \(Y\) after fixing \(A\).

To prove identification for the enlarged action set, take any model
on the stated graph and any finite alphabets with positive observed
law. Every latent input used by \(Y\) is private to \(Y\),
since a shared input would create a bidirected edge incident with
\(Y\), absent from the graph. Its input tuple is therefore
independent of all inputs determining \(A,M\). Its only observed
parent is \(M\). Define
\(k(y;m)=\mathbb P(f_Y(m,U_{\Lambda_Y})=y)\).
Conditioning on the positive event \(M=m\) leaves the independent
\(Y\) input law unchanged, so \(p(y\mid m)=k(y;m)\).
Replacing both equations \(A,M\) makes the response equation
exactly \(f_Y(m,U_{\Lambda_Y})\), and therefore
\(P_{a,m}(y)=k(y;m)=p(y\mid m)\) for all assignments.
This proves graph-uniform identification of the enlarged query.
In the displayed models the kernel's response-one values are
\(1/4\) at \(m=0\) and \(3/4\) at \(m=1\).

The smaller forest of the earlier hedge contains \(M\), which
is now an action vertex. It fails the hedge condition for the
enlarged action set. Thus the correctly stated criterion keeps
all original actions in both the smaller-forest exclusion and the
post-action ancestor test. Existence of an obstruction for an
arbitrarily chosen subset of those actions does not suffice.
\end{proof}
\end{annexitem}

\begin{annexitem}{annex:p16.tian-component-invariant}{thm:p16-tian-component-invariant}{The exact district inclusions and parameter-preserving support invariant}
\begin{proof}
The ancestor set \(A=\operatorname{an}_{G[T]}(C)\) contains
\(C\) and is ancestral within \(T\). If \(A=C\), the first
branch applies, including the case \(C=T\). If it does not
apply and \(A=T\), the second branch has precisely
\(C\subsetneq T=\operatorname{an}_{G[T]}(C)\).
Otherwise \(C\subsetneq A\subsetneq T\), so these branches
exhaust all cases in the stated order.

In the recursive case the ancestral summation theorem gives
\(q_A=\Sigma_{T\setminus A}q_T=Q[A]\). Every outside
assignment is held fixed in this sum. The removed coordinates
\(T\setminus A\) become additional fixed arguments of
\(Q[A]\), but its value is independent of them because they
are not parents of vertices in the ancestral set \(A\).
At each retained assignment, this marginal is a nonempty finite
sum of positive values, hence positive; as an actual intervention
kernel it is normalized on \(S_A\).

The induced graph on \(C\) remains bidirected-connected inside
\(G[A]\), so its vertices belong to one unique maximal district
\(T'\). Thus
\(C\subseteq T'\subseteq A\subsetneq T\).
Neither of the first two inclusions has to be strict. Nonemptiness
and connectivity of \(T'\) follow from district membership.
Every prefix marginal of the positive \(q_A\) is positive,
so every ratio used in \(\operatorname{Ext}^A_{T'}(q_A)\)
has a positive denominator. The district recovery theorem identifies
the resulting positive product with the actual normalized kernel
\(Q[T']\). Its outside arguments are all of \(N\setminus T'\),
including newly fixed coordinates in \(A\setminus T'\); those
values are retained throughout the extraction rather than conditioned
away. This proves the kernel invariant for the child. Induction
along recursive calls proves it for every reached call.
\end{proof}
\end{annexitem}

\begin{annexitem}{annex:p16.tian-component-soundness}{thm:p16-tian-component-soundness}{Finite component recursion and correctness of every return}
\begin{proof}
Only branch 3 recurses, and its child has
\(|T'|<|T|\) while retaining the same nonempty \(C\).
Thus the nonnegative integer \(|T|-|C|\) strictly decreases
at each recursive edge. There is one child at each such edge,
so the recursion terminates after at most \(|T|-|C|+1\) calls.
All finite graph operations, sums and ratios are defined by the
preceding invariant. A failure only shortens a returning path.

Induct on \(|T|-|C|\) for soundness. At branch 1, \(C\) is
ancestral within \(T\), so the already proved ancestral identity
gives \(\Sigma_{T\setminus C}q_T=Q[C]\), with all outside
parameters retained. Branch 2 gives no expression. At branch 3,
the child is smaller and has input \(Q[T']\). If it succeeds,
induction gives its target as \(Q[C]\), the very same global
kernel required by the parent. The child does not change the
meaning or fixed arguments of that target. This proves every
successful return without an assumption that a failure is meaningful
for an arbitrary externally restricted model class.

Every produced expression uses a finite ancestral sum or a finite
product of positive prefix ratios, possibly substituted into the
smaller expression. Therefore it is a finite expression of the
asserted kind in the original input kernel. For the distinct
do-calculus route, view \(q_T\) as the positive law of the
restricted model after fixing \(N\setminus T\). The source
ancestral reduction identity, proved through Rule 3 in
\cref{thm:p16-id-do-utility}, licenses the ancestral restriction;
ordinary summation computes its input law. The district extraction
uses that theorem's complete source Lemma 4 derivation. At branch 1,
the same ancestral restriction identifies \(Q[C]\) with the
computed marginal. Inserting child derivations into parent derivations
along the finite recursion gives the claimed do-calculus route.
Its local-to-original graph and support bridge is already proved
in \cref{thm:p16-id-do-licenses}.
\end{proof}
\end{annexitem}

\begin{annexitem}{annex:p16.tian-outer-soundness}{thm:p16-tian-outer-soundness}{The full outer kernel calculation and all free parameters}
\begin{proof}
The empty response is a unit-mass empty tuple under every intervention,
so the explicit first return is correct. Otherwise \(Y\subseteq D\)
and \(D\) is nonempty. Each maximal district \(D_i\) of
\(G[D]\) is bidirected-connected in the full graph and therefore
lies in a unique full-graph maximal district \(T_i\).
Extraction from the positive observed law gives
\(q_{T_i}=Q[T_i]>0\) at every full assignment by
\cref{thm:p16-latent-district-recovery}. Thus every initial component
call satisfies its interface, and the component invariant proves
positivity and valid outside parameters in all its descendants.
There are finitely many districts, and each call terminates, so
the outer procedure terminates.

If every child succeeds, its value is exactly \(e_i=Q[D_i]\).
The proved district product gives
\(\prod_i e_i=Q[D]\). The target reduction theorem then gives
\[
 \Sigma_{D\setminus Y}\prod_i e_i
   =\Sigma_{D\setminus Y}Q[D]
   =P_x(y).
\]
All coordinates are interpreted through the same full assignment.
Coordinates of \(D\setminus Y\), even if they were outside
arguments of another child's factor, are bound by this parent sum.
The original \(X\) coordinates are never summed because
\(D\subseteq N\setminus X\).

To verify the remaining parameter assertion directly, an observed
parent of any vertex in \(D\), if outside \(X\), has a
directed path to \(Y\) in \(G[N\setminus X]\). It must
therefore itself belong to \(D\). Hence every outside parent
of \(D\) belongs to \(X\). The kernel formula for \(Q[D]\)
uses outside values only at those parents, so it is independent
of \(N\setminus(D\cup X)\); the final marginal inherits that
independence. A reference assignment on those coordinates is
therefore immaterial, and no positive observational event there
has been assumed as a condition for the structural intervention.

For the do-calculus derivation, original nonancestors after \(X\)
can be fixed to arbitrary values by the proved additional-action
identity, the source Lemma 6 route. This reduces the target to
\(\Sigma_{D\setminus Y}Q[D]\). The source Lemma 4 route
supplies its district product. Each full-graph district extraction
and each successful component call has the derivation proved
above. Substitution gives a uniform finite derivation ending in
the observed-law formula. All conditional denominators remain
positive by the invariant. Finally, on a specified singleton family,
the previous elimination/restoration theorem transports a successful
reduced answer to the same original response kernel.
\end{proof}
\end{annexitem}

\begin{annexitem}{annex:p16.tian-failure-hedge}{thm:p16-tian-failure-hedge}{Explicit common-root forests and the original action intersection}
\begin{proof}
The selected outer component \(C=D_i\) never changes during
its recursion. It is a maximal district of the original post-action
ancestor graph \(G[D]\), and is disjoint from the original
action set \(X\). At failure, the first branch did not apply,
so the invariant gives
\(C\subsetneq T=\operatorname{an}_{G[T]}(C)\), with both
induced graphs bidirected-connected.

Suppose \(T\cap X=\varnothing\). Every vertex in \(T\)
has a directed path inside \(T\) to \(C\). Each vertex
of \(C\subseteq D\) has a directed path to \(Y\) avoiding
\(X\). Their concatenation shows that all of \(T\) lies
in \(D\). Since \(T\) is bidirected-connected there and
contains the maximal district \(C=D_i\), maximality forces
\(T\subseteq C\), contradicting its proper inclusion. Thus
\(T\cap X\ne\varnothing\). This proves the original action
intersection from the actual recursion; the mere fact that the
initial containing district was larger than \(D_i\) would not
by itself suffice for the final failed set.

Let \(K\) be the directed sinks of \(G[C]\). Every vertex
of the finite nonempty DAG \(G[C]\) reaches a sink by repeatedly
following an outgoing edge, so \(K\ne\varnothing\) and every
vertex in \(C\) has a path to \(K\) within \(C\).
Choose a first edge on a shortest such path at each \(c\in C\setminus K\),
and give \(K\) no outgoing edge. At every \(t\in T\setminus C\),
choose a first edge on a shortest path to \(C\) inside \(T\).
Distances decrease until \(C\) is reached, then toward \(K\).
The selected directed graphs on \(T\) and \(C\) therefore
have at most one child per vertex and exactly the same sinks \(K\).
In particular the construction is valid even if a sink of the
unpruned \(G[C]\) originally had a child outside \(C\).

Choose a bidirected spanning tree on \(C\) and extend it
one new adjacent vertex and edge at a time to a spanning tree
on \(T\). Connectivity guarantees each extension, as in the
earlier hedge extraction proof. The resulting edge subgraphs are
nested C-forests with common roots \(K\). Their smaller vertex
set misses \(X\), their larger set meets \(X\), and
\(K\subseteq C\subseteq D=\operatorname{an}_{G[N\setminus X]}(Y)\).
They form a hedge for the original query. Retaining every containing
frame preserves the original graph, query, fixed arguments and
selected district associated with this construction.

For a standalone component call beginning at \(T_0\), the
failed set \(T\) is a subset of \(T_0\) and strictly contains
\(C\). Thus it meets the standalone action set \(T_0\setminus C\),
while \(C\) misses that set. The same roots \(K\) belong
to the standalone response \(C\) itself. The same selected
forests are therefore a hedge in \(G[T_0]\) for that standalone
complementary-intervention query. No relation between unpruned
root sets was needed in either argument.
\end{proof}
\end{annexitem}

\begin{annexitem}{annex:p16.tian-completeness}{thm:p16-tian-completeness}{Completeness of both procedures and equality with the ID decision}
\begin{proof}
The outer procedure terminates and all its successes are sound
for every compatible positive law on every finite nonempty alphabet
family. Its graph tests, like those of ID, do not use numerical
values or alphabet cardinalities. Thus success supplies one uniform
finite expression and establishes graph-uniform identification.

If it fails, the previous theorem extracts a hedge for its original
action and response sets. The positive four-valued routing construction
produces two whole-graph models, on an allowed finite alphabet family,
with equal full observed laws and different original effects. That
one family disproves graph-uniform identification. This constructs
countermodels on the original graph; it does not assume that a
nonidentified child factor stays nonidentified under every possible
outer sum, or attempt to embed an arbitrary supplied child kernel
into a prespecified full observed law.

Termination, success soundness and this negative direction prove
outer completeness. Both the outer procedure and ID therefore
succeed exactly on the uniformly identified queries. Their graph
decisions agree; whenever successful, their expressions both equal
the same structural intervention kernel in each compatible model.
Their products and sums may be arranged differently, so no syntactic
identity is asserted. The graph-uniform all-effects theorem now applies
to identify as well through this equality of graph decisions.

For the standalone component question, regard \(G[T_0]\) as
the whole observed graph and its positive law as the supplied input.
Equivalently, any outside parameters have been fixed in advance.
Success soundness gives its complementary-intervention kernel on
\(C\). A failed call extracts a hedge for action
\(T_0\setminus C\), response \(C\), with roots already
in that response. Apply the positive binary forest witness with
empty external routing, and extend by independent positive coordinates
in \(T_0\) outside the failed forest, retaining all declared
edges as in the established embedding proof. It gives equal positive
input laws and different standalone target kernels. The binary
alphabet lift supplies the same conclusion for every prescribed
family on \(T_0\) with at least two values per vertex. It also
disproves the graph-uniform quantifier, which includes such families.
Hence c-identify is complete for the stated standalone question.
This theorem quantifies over possible input laws of that class;
it makes no claim of a countermodel with any particular fixed
numerical kernel chosen by an external specification.
\end{proof}
\end{annexitem}

\begin{annexitem}{annex:p16.tian-worked}{thm:p16-tian-worked}{All calls, selected districts and numerical kernels in the examples}
\begin{proof}
\paragraph{The selected district equals the target C.}
For the first graph take independent fair bits \(U,V\) and
independent bits \(E_Z,E_Y\), each of one-probability \(1/4\),
all four mutually independent. Use
\[
 X=U,\qquad Z=U\oplus V\oplus E_Z,\qquad
 Y=X\oplus V\oplus E_Y.
\]
The common inputs \(U,V\) have child sets \(\{X,Z\}\)
and \(\{Z,Y\}\), and the other inputs are private. Thus the
declared graph is exactly the first stated graph. Write
\(q(0)=3/4,q(1)=1/4\). For a proposed assignment,
\[
 p(x,z,y)=\tfrac14\sum_{v=0}^1
      q(z\oplus x\oplus v)q(y\oplus x\oplus v).
\]
If \(z=y\), the sum is \((3/4)^2+(1/4)^2=10/16\),
and the mass is \(5/32\). If \(z\ne y\), it is
\(2(3/4)(1/4)=6/16\), giving \(3/32\).
There are four assignments of each type, so these positive masses
sum to one.

The initial connected set is \(T=\{X,Z,Y\}\), its kernel
is \(p\), and its target is \(C=\{Y\}\). The only directed
ancestor of \(Y\) other than itself is \(X\), so
\(A=\{X,Y\}\). Both the base and failure branch are skipped.
The ancestral marginal is
\[
 q_A(x,y;z)=\sum_{z'}p(x,z',y)=(5+3)/32=1/4,
\]
independent of its displayed fixed \(z\) argument. The graph
on \(A\) has two singleton districts, so the district containing
\(C\) is exactly \(T'=C=\{Y\}\). In topological order
\(X,Y\), its extracted kernel is
\[
 q_{T'}(y;x,z)=\frac{q_A(x,y)}{\sum_{y'}q_A(x,y')}
             =\frac{1/4}{1/2}=1/2.
\]
The recursive call now has \(C=T'=\{Y\}\); its ancestor
set equals \(C\), so branch 1 returns this kernel with an
empty summation set. The fixed \(x,z\) parameters remain
typed, although this particular result ignores both. Directly
replacing \(X,Z\) leaves \(Y=x\oplus V\oplus E_Y\),
which is fair, confirming the calculated kernel.

\paragraph{The selected district equals the ancestor set A.}
For the second graph use independent \(U,E,W\), where \(U,W\)
are fair and \(\mathbb P(E=1)=1/4\). Use \(X=U,Z=W\)
and either \(Y=X\oplus U\oplus E\) or \(Y=E\).
Declare \(U\) shared by \(X,Y\), and add an ignored
independent constant input shared by \(Y,Z\) to supply its
declared bidirected edge. Both equations may ignore that input.
The common observed law is \(p(x,y,z)=\frac14 q(y)\),
giving the stated positive masses \(3/16\) and \(1/16\).

Again \(C=\{Y\}\), \(T=\{X,Z,Y\}\), and
\(A=\{X,Y\}\). Now the graph on \(A\) contains
\(X\leftrightarrow Y\), so its only district is \(T'=A\).
The ancestral marginal is
\(q_A(x,y;z)=\sum_{z'}p(x,y,z')=\frac12q(y)\).
Extracting the only district multiplies all prefix conditionals
and returns \(q_{T'}=q_A\) itself. Its first marginal is
\(q_X(x)=1/2\); all conditional denominators are positive.
The new call is \((C=\{Y\},T'=\{X,Y\},q_A)\).
Its ancestor set is all of \(T'\), because \(X\to Y\),
and is larger than \(C\). Thus its first branch does not apply
and its second branch returns the claimed failure. The raw bow
and singleton \(Y\) give its hedge. In these two models the
response-one effect under fixing \(X\) is respectively \(1/2\)
and \(1/4\), regardless of any additional action on \(Z\).
This example shows why equality \(T'=A\) is a legitimate
recursive case even though the next call can fail.

\paragraph{The complete front-door return.}
For the third graph and target \(X=\{A\},Y=\{C\}\),
the post-action ancestor set is \(D=\{B,C\}\), whose districts
are \(\{B\},\{C\}\). The full-graph districts are
\(\{B\},\{A,C\}\). Their extracted kernels, in order
\(A,B,C\), are
\[
 q_B(b;a,c)=p(b\mid a),\qquad
 q_{AC}(a,c;b)=p(a)p(c\mid a,b).
\]
The first component call has equal target and input set \(\{B\}\)
and returns \(q_B\) at branch 1. The second has target
\(\{C\}\), input set \(\{A,C\}\), and kernel \(q_{AC}\).
The induced graph on \(A,C\) has no directed edge, so its
target ancestor set is just \(\{C\}\); branch 1 returns
\[
 q_C(c;a,b)=\sum_{a'}p(a')p(c\mid a',b).
\]
Its formerly random \(A\) coordinate is now a displayed fixed
parameter which the marginal ignores. The dummy \(a'\) is
distinct from the original action assignment \(a\). The outer
step multiplies \(q_B q_C\) and sums \(B\), giving exactly
the stated formula, with the original \(a\) retained.

For the positive worked model in \cref{thm:p16-id-worked-success},
the already calculated inner response-one kernel is \(1/2\)
at \(b=0\) and \(7/8\) at \(b=1\). The mediator's
conditional probabilities are \((3/4,1/4)\) at \(a=0\)
and \((1/4,3/4)\) at \(a=1\). Thus the outer sums are
\[
 \tfrac34\tfrac12+\tfrac14\tfrac78=\tfrac{19}{32},\qquad
 \tfrac14\tfrac12+\tfrac34\tfrac78=\tfrac{25}{32}.
\]
That model's full table is strictly positive; its relevant prefix
marginals are \(p(a)=1/2\) and
\(p(a,b)=\frac12p(b\mid a)>0\). Thus every extraction and
conditional in this trace has its required positive denominator.

\paragraph{Ancestor equality does not fix the unpruned roots.}
In the fourth graph the induced graph on \(C=\{a,b\}\)
has its bidirected edge and no directed edge, so both \(a,b\)
are sinks. In \(T=\{a,t,b\}\), the directed path
\(a\to t\to b\) leaves only \(b\) a sink. Every vertex
in \(T\) is an ancestor of a member of \(C\), so
\(T=\operatorname{an}_{G[T]}(C)\) nevertheless holds.
Both sets are bidirected-connected, and the component call fails
because \(C\subsetneq T\). Choose the common root set
\(K=\{a,b\}\). The prescribed construction removes the
outgoing edge \(a\to t\) at the chosen root \(a\), keeps
\(t\to b\) to route the remaining vertex into \(C\),
and keeps the bidirected tree edges \(a\leftrightarrow b\),
\(a\leftrightarrow t\). Its restriction to \(C\) has
no directed edge. Both resulting forests now have exactly the
roots \(K\), proving each assertion without the invalid
unpruned-root equality.
\end{proof}
\end{annexitem}
